\section{A company program begins with the state it can justify}
\label{sec:company-source-state}

P's return example has an opening book, identified lots, dated events and
explicit tax premises. That information allows the reference calculation to
reach an annual result. An actual company's published statements usually
offer a different starting point: amounts reported for a particular group
and period, together with policies and notes that leave some underlying
contracts and allocations unknown. The language must accommodate both
starting points. An author should be able to select methods for a complete
conditional example and also describe a source-constrained partial state
whose missing facts still prevent those methods from running.

\subsection{Selecting methods for a complete retained company}

For P, the finite composed-company declaration names the source packet and
opening, legal-book context, reporting period, return component and all
thirteen events in their original order. Each event selects a declared
capability. Return settlement and recovered-lot sale use the native return
method; the two tax closes use the native tax method; the other transactions
use the previously checked corporate methods. Annual reporting and carry
have their own explicit selections. The caller provides the available
implementations and their identities. A capability name written in a
program cannot import arbitrary code, and an unavailable capability cannot
silently select a different method.

Matched text and typed records express the same declaration. The relevant
draft is \path{RetainedCompanyProgramDraft/1}; its example and evaluator
are identified at the end of this section. Binding the program to the
retained case means comparing its full source, opening, event schedule and
method selections, rather than accepting a familiar company name as proof
of identity. An author cannot rename P as Apple and retain the old source
facts. Nor does a correctly hashed result establish its own correctness:
the reference verifier recomputes the proposals, annual report and carry.

Review findings remain separate from immutable source facts. If a required
tax finding is unknown, the first seven valid events remain available,
the six subsequent events remain unexecuted, and no complete annual report
or next opening is returned. A negative finding instead records that the
selected treatment is unsupported. Editing a finding changes the authored
program's identity; it does not edit the source packet. With the retained
findings supplied, the composed calculation reproduces P's profit of 49.5,
cash of 491.5, equity of 671.5 and operating cash flow of 101.5. This
completes the finite declaration for the case while leaving broader company
composition and method admission to their own specifications.

\subsection{What the disclosed balance sheet leaves unresolved}

Consider the retained Apple balance sheet at 28 September 2024. It reports
cash and cash equivalents of 29,943, inventory of 7,286 and net property,
plant and equipment of 45,680, in millions of US dollars
\citep[p.~31]{apple2024}. These are useful observations. They do not, however,
identify the unrestricted cash available to each legal entity, inventory
quantities and lot costs, or the original costs, accumulated depreciation,
service dates and remaining lives of asset cohorts. Supplying one invented
unit cost or useful life would choose an economic history that the
disclosure has not established.

The input-readiness reference therefore preserves the original statements
and asks what each selected execution method still requires. It keeps all
139 package-inventory entries, the 102 annual and 100 interim readings,
and the two collections of 27 opening or comparative balances. Fifteen
dependencies connect those records and the original note obligations to
company inputs, affected operations, reports and annual carry. Missing
claims, asset cohorts, tax bases and dated events remain visible beside the
known displayed amounts.

There are also restrictions in the selected methods themselves. The
operations draft requires a full calendar year, a fixed account registry,
homogeneous inventory cost, specified asset cohorts and a simple tax
treatment. Apple's retained consolidated disclosures do not satisfy those
requirements merely because their balance sheet balances. P's composed
program has a different restriction: it binds exactly the retained
synthetic case. The readiness result keeps missing company state and
unsupported method scope as separate reasons for refusing execution.
That refusal concerns the selected profiles; it does not establish that
Apple cannot be represented by a broader model.

Known arithmetic can still be retained without pretending to solve those
problems. The displayed current and noncurrent term-debt captions sum to
\(10{,}912+85{,}750=96{,}662\) at the annual date. This is a sum of two
reported captions, not total financing obligations or an allocation to
individual instruments. The source's dash for annual debt-issuance cash
flow remains a dash. Unknown original-byte publication and ingestion clocks
limit historical as-of admission, but they do not prevent reading the
named retained source. The unresolved actual-company history and the earlier
conditional subsystem calculations remain separate, as before.

\subsection{Fiscal dates belong to the source}

A partial-state declaration can carry the reported measurements without
manufacturing an executable opening. The draft
\path{SourceConstrainedCompanyStateDraft/1} binds the exact readiness request
and result, declares its fiscal intervals and names each selected observation.
It includes all 54 current annual and interim balance captions. Each binding
keeps the original measurement date, consolidated perimeter, basis, units,
precision and source-admission status. Prior comparisons and the remaining
readings stay available in the referenced readiness result; they have not
all become executable state constraints.

The retained annual source covers 1 October 2023 through 28 September 2024;
the following interim source covers 29 September through 28 December 2024
\citep{apple2024,apple2025q1}. Using the half-open interval convention already
introduced in Chapter~\ref{ch:history}, these become
\begin{equation}
\label{eq:source-fiscal-periods}
  [2023\text{-}10\text{-}01,\,2024\text{-}09\text{-}29),
  \qquad
  [2024\text{-}09\text{-}29,\,2024\text{-}12\text{-}29).
\end{equation}
The exclusive end is the day after the source's inclusive closing date.
The first interval has 364 days and the second 91; they meet without a
gap or overlap. A balance-sheet observation dated 28 September remains at
that date, rather than moving to the exclusive boundary on 29 September.

Those date calculations specify reporting coverage. They do not select an
interest day-count convention, a depreciation fraction or tax-year law.
The generic structural validator admits explicit non-calendar periods;
the retained-source binder separately rejects a period that disagrees with
the original observation. Earlier fiscal-closing examples already permit
short periods. The addition here makes fiscal intervals directly authorable
alongside a partial source state, while preserving the stricter scope of
the existing execution methods.

\subsection{A rounding assumption constrains a measurement}

The default interpretation of a displayed amount is deliberately modest:
the source displays this value, while its unrounded underlying measurement
is unresolved. Storing the displayed integer as an exact rational preserves
the token; it does not establish exact equality with the company's
underlying balance. Even an interval requires a rounding convention that
the retained precision annotation does not by itself establish.

An author may supply a separate rounding hypothesis. Suppose the reporting
quantum is \(q>0\), the displayed value \(d\) is an integer multiple of
\(q\), and the rule rounds to the nearest multiple with ties away from
zero. Let \(x\) denote the underlying value of the same reported
measurement: it has the same date, perimeter, unit and basis, not a newly
inferred legal-account meaning. The midpoints adjacent to \(d\) are
\(d-q/2\) and \(d+q/2\). At a midpoint the value with the larger absolute
magnitude wins. It follows that the values producing \(d\) are
\begin{equation}
\label{eq:source-rounding-preimage}
  x\in
  \begin{cases}
    [d-q/2,\,d+q/2), & d>0,\\
    (d-q/2,\,d+q/2], & d<0,\\
    (-q/2,\,q/2), & d=0.
  \end{cases}
\end{equation}
For positive \(d\), the lower tie rounds up to \(d\), while the upper tie
rounds away to \(d+q\). For negative \(d\), the endpoint inclusion
reverses. Neither midpoint rounds to zero, which explains the two open
endpoints in the last case.

Under that explicit hypothesis with \(q=1\) million dollars, the annual
cash-and-equivalents caption of 29,943 permits
\([29{,}942.5,29{,}943.5)\). The displayed retained deficit of
\(-19{,}154\) permits \((-19{,}154.5,-19{,}153.5]\)
\citep[p.~31]{apple2024}. These are conditional deductions, not evidence
that Apple uses the proposed tie-breaking rule. The assumption has its own
identity and provenance, and it remains distinct from accepted source facts.
Without it, the declaration retains the displayed value and leaves the
underlying measurement unresolved.

The intervals neither allocate cash to legal entities nor split inventory
into lots. They are not assumed independent, and their Cartesian product
is not asserted to be a feasible company history. Every missing-state
dependency therefore survives the rounding calculation. The reference
returns a partial description, with joint feasibility and actual history
unknown, and supplies no executable opening. This lets the language retain
what a source says while making the next evidence or modelling choice
precise enough to review.

The closed typed records, named text examples and bounded reference checks
are under \path{data/accounting-master/}:
\begin{quote}\small
\path{company-composed-profile-20261024}\\
\path{real-company-input-readiness-20261024}\\
\path{source-constrained-state-20261024}
\end{quote}
Their checks cover source identity, complete dependency retention, matching
text and typed interpretations, fiscal boundaries and signed rounding ties.
They reject changed source dates, legal-book relabelling, invented exact
balances and a forged executable result. The next section extends the source
description beyond current balance captions while preserving the difference
between displayed relations and latent company state. General source binders, admissible
state reconstruction, broader accounting methods and independent acceptance
of the governing design remain unfinished.

\section{Reading the statements together without inventing history}
\label{sec:statement-context-relations}

The balance sheet gives a dated description of the company. The income,
cash-flow and equity statements explain different aspects of the movement
between dates. Reading them together requires more than adding amounts with
similar labels. An equity balance is measured at an instant, profit accrues
over a period, a cash-flow adjustment can concern a different population of
items from the balance-sheet caption, and a subtotal calculated by the
analyst supplies no independent observation against which to check itself.

The retained Apple statement networks already record those relationships
and unresolved differences. The complete context draft makes their roles
explicit in the language. Its combined view preserves 256 source observations
and all 143 original relation records, including their original explanations
and diagnostic classifications. It gives 135 displayed differences, six
derived subtotals and two calculations that remain unresolved because an
operand is a reported dash. This accounts for the full retained network;
it does not establish the missing underlying company history.

\subsection{A reporting period can contain several balance dates}

Apple's annual equity statement is a report for fiscal 2024, but its opening
equity is measured at 30 September 2023 and its closing equity at
28 September 2024 \citep[p.~32]{apple2024}. The portable observation retains
the annual report-period string. The original pointed reading separately
records the measurement date. The adapter follows that pointer and retains
both meanings: the date of the stock and the period of the statement in
which it appears. It does not overwrite one with the other.

The same distinction matters for cash. An opening cash amount is a stock;
operating cash flow is a movement during the period. A cash bridge joins
the two endpoints and that intervening movement. A snapshot subtotal instead
compares amounts measured at the same instant. The language therefore
declares the relation's purpose as well as its operands. Same-period income
comparisons, same-date equity comparisons, stock-flow bridges, net-asset
changes and explanatory differences have different context rules.

For example, the displayed annual change in net assets is compared with the
displayed change in equity, using both sets of reported endpoints. With
\(A_t\), \(L_t\) and \(E_t\) denoting assets, liabilities and equity in
millions of dollars at the closing date, and subscript \(0\) denoting
the opening date, the comparison is
\begin{align}
\label{eq:statement-context-equity-change}
 (A_t-A_0)-(L_t-L_0)
 &= (364{,}980-352{,}583)-(308{,}030-290{,}437)\notag\\
 &=12{,}397-17{,}593=-5{,}196\notag\\
 &=56{,}950-62{,}146=E_t-E_0.
\end{align}
All operands come from the retained balance-sheet readings
\citep[p.~31]{apple2024}. Recomputing the right side from the equity
endpoints gives it a distinct observed basis. The corresponding interim
comparison gives 9,808. These displayed reconciliations retain the original
precision and consolidated scope; they do not identify the transactions
that produced the changes.

An annual-to-interim endpoint comparison also retains both source contexts.
The annual closing equity and the next quarter's opening equity are
reported at the same date, even though they appear in different reports.
Their equality can be compared in the combined named-source view. An
interim-only view cannot use the unselected annual observation. Its five
cross-period comparisons therefore remain unresolved, with their original
relation records still visible. No later source is silently supplied to
complete them, and the annual and interim bases are not generally declared
interchangeable.

\subsection{A subtotal needs an independent target to be a check}

The source audit exposed an evidentiary weakness in the earlier annual
relationship ledger. Three analyst-calculated totals---other equity
movements, operating reconciliation adjustments and operating asset/liability
adjustments---were stored as fixed numeric targets and labelled displayed
pass/fail checks. Reproducing those totals checks the calculation against
its saved value. It does not supply an independent disclosed total or an
additional accounting identity. The interim ledger already treated the
corresponding calculations as explanatory subtotals.

The new interpretation classifies all six as derived quantities and returns
no independent residual for them. Their original constants, coefficients
and classifications remain in the retained records, so the reason for the
change can be examined. The source amounts and earlier arithmetic are not
rewritten. The correction concerns the strength of the evidence attributed
to a calculated subtotal.

Other differences have a different purpose. The annual opening cash-flow
population exceeds balance-sheet cash by 772. The annual payable comparison
leaves 329, and interim tax expense less cash tax paid leaves
\(-12{,}397\). Those differences remain explanations to investigate, not
amounts to post into a balancing account. Even the zero vendor-receivable
difference does not identify a collection schedule, gross transactions or
their dates. The full-context result retains each relation's populations
and declared diagnostic purpose beside its arithmetic.

Two annual calculations cannot be completed under the new baseline because
the debt-issuance operand is a dash. The earlier displayed-total calculation
used an explicit no-contribution convention. That convention and its result
remain in the historical evidence, while the new source interpretation
preserves the dash and leaves the financing total and term-debt diagnostic
unresolved. Using such a convention again would require its own explicit
premise. Similarly, five interim relations retain the annual concepts absent
from the interim table. Summing the rows that are actually listed does not
establish that an omitted class of transaction was zero.

\subsection{Local compatibility is not a company history}

The optional rounding intervals in Section~\ref{sec:company-source-state}
make it possible to ask a further, conditional question: could an individual
measurement identity hold for values inside those intervals? This requires
two separate premises. One specifies the rounding rule; the other asserts
that the proposed identity concerns the same underlying measurements.
A zero displayed residual alone supplies neither premise.

Let a proposed relation have residual \(r=\sum_i a_i x_i\), including
the target with coefficient \(-1\). Before computing an interval for
\(r\), collect repeated occurrences of the same observation and add their
coefficients. Thus \(x-x=0\) exactly, rather than the difference of two
independently varying copies of \(x\). For a positive coefficient, the
smallest contribution uses the lower endpoint and the largest uses the
upper. For a negative coefficient those choices reverse. Summing the
contributions gives the residual endpoints; an endpoint belongs to the
result only if every nonzero contribution can attain its relevant endpoint.

For the annual current-assets relation, six positive displayed components
sum to the displayed total of 152,987. Under the proposed nearest-rounding
rule with quantum one million, write the component errors as
\(e_1,\ldots,e_6\) and the total's error as \(e_T\). Every one of these
positive displayed measurements has error in \([-1/2,1/2)\). The residual
of the underlying measurements is therefore
\begin{equation}
\label{eq:statement-context-residual-range}
  r=\sum_{i=1}^{6}e_i-e_T\in(-7/2,\,7/2).
\end{equation}
The lower extreme would require the total's excluded upper endpoint, and
the upper extreme would require the components' excluded upper endpoints.
Zero lies inside this range. Under the separately proposed identity, that
establishes local compatibility with the intervals. It leaves the rest of
the statement network untested.

A simple counterexample shows why this distinction is necessary. If
\(x\in[0,1]\), the proposed equality \(x=0\) is locally compatible,
and so is \(x=1\). No single value satisfies both. A collection of local
interval tests is therefore not a witness for a joint feasible history.
The current reference reports joint feasibility as uncomputed and actual
history as unknown. It creates neither legal books nor the claims, lots,
contracts and behavioral assumptions required to continue the company.

The complete observation and relation declarations have matched typed and
named-text forms in \path{RetainedStatementContextDraft/1}. Each operand
has its source identity, date or interval, unit and explicit alias. Relation
roles, signs, targets, source omissions and retained dependencies are part of
the declaration. Source binding rejects changed dates, units, populations,
aliases and roles; result verification recomputes the complete description.
The reference preserves per-share values and weighted-average share counts
alongside monetary amounts without summing unlike units.

The checked records are under \path{data/accounting-master/}:
\begin{quote}\small
\path{source-relations-20261024}\\
\path{statement-context-coverage-20261024}
\end{quote}
The former contains conditional interval diagnostics for the balance
relations; the latter supplies complete displayed-statement contexts and
explicit unresolved results. Together they make the source layer more
precise. General source-to-state mappings, shared method requirements,
admissible latent histories and independent acceptance remain subsequent
design obligations.

